\documentclass[11pt]{amsart}

\usepackage[margin=1in]{geometry}
\usepackage{amsmath,amssymb,amsthm}
\usepackage{mathtools}
\usepackage{booktabs}
\usepackage{array}
\usepackage[colorlinks=true,linkcolor=blue,citecolor=blue,urlcolor=blue]{hyperref}
\usepackage{enumitem}
\usepackage{xcolor}

% --- theorem environments ---
\theoremstyle{plain}
\newtheorem{theorem}{Theorem}[section]
\newtheorem{proposition}[theorem]{Proposition}
\newtheorem{lemma}[theorem]{Lemma}
\newtheorem{corollary}[theorem]{Corollary}
\newtheorem{conjecture}[theorem]{Conjecture}

\theoremstyle{definition}
\newtheorem{definition}[theorem]{Definition}
\newtheorem{example}[theorem]{Example}
\newtheorem{observation}[theorem]{Observation}

\theoremstyle{remark}
\newtheorem{remark}[theorem]{Remark}

% --- macros ---
\newcommand{\W}{W}
\newcommand{\Wq}{W(q)}
\DeclareMathOperator{\inv}{inv}
\DeclareMathOperator{\Var}{Var}
\newcommand{\TODO}[1]{\textcolor{red}{\textbf{[TODO: #1]}}}

\title[Log-concavity of Bruhat intervals in Weyl groups]
{Log-concavity of Bruhat intervals in Weyl groups:\\
extended verification, an extremal conjecture, Mahonian asymptotics,\\
and equality cases}

\author{Nikol Panayotova Savova}
\address{University of Oxford}
\email{nikol.p.savova@gmail.com}

\author{Sihao Huang}
\address{Independent Researcher}
\email{sihao.c.huang@gmail.com}

\date{\today}

\begin{document}

\begin{abstract}
Brenti conjectured that the rank sequence of every Bruhat interval in a
finite Weyl group is log-concave; the statement is known to fail for the
non-crystallographic Coxeter group $H_3$. We report an extended
computational verification of the conjecture, exhaustively checking
$1{,}079{,}490{,}991$ Bruhat intervals across the Weyl groups of types
$A_2$--$A_7$, $B_2$--$B_6$, $D_4$--$D_6$, $E_6$, $F_4$, and $G_2$ in exact
integer arithmetic, with zero violations found; this extends the previously
recorded verification frontier by roughly an order of magnitude and closes
several cases ($A_6$, $A_7$, the full range of $B_5$, $B_6$, $D_6$, $E_6$)
left open in Brenti's own list. We complement this with non-exhaustive evidence
in higher rank: near-top slab sweeps in $A_8$--$A_{10}$, $D_7$--$D_8$,
$E_7$, and $124{,}944$ sampled and seeded intervals in $B_7$, $B_8$, $D_7$,
$E_7$. The data suggest an extremal principle: in every irreducible,
simply-laced Weyl group the global minimum of the log-concavity ratio over
all Bruhat intervals is attained by the full interval $[e,w_0]$; we show
both hypotheses --- irreducibility and the simply-laced restriction --- are
necessary, and record a proved instance of the principle restricted to
rationally smooth intervals. In type $A$ this minimum is the central ratio
of the Mahonian distribution; we prove an unconditional local limit theorem
pinning its value near the center for every $m$, determine the global
minimum exactly, in integer arithmetic, for every $4\le m\le560$, and prove the
sharp asymptotic $\sigma_m^2(r_m-1) = 1-\tfrac{27}{25}m^{-1}+O(m^{-2})$
conditional on a single explicitly stated core lemma on the exponentially
tilted Mahonian distribution --- needed only for $m\ge561$, and itself
reduced, with explicit constants throughout, to four explicitly stated open
cumulant statements. We also
record a classification of the observed equality cases as rank-two dihedral
parabolic patterns of order $m\in\{4,6\}$, together with a candidate
explanation, via the crystallographic restriction, for why Weyl groups
escape the failure Brenti records for the non-crystallographic group $H_3$,
whose counterexample interval sits on a dihedral core of order $m=5$.
\end{abstract}

\subjclass[2020]{05E15, 20F55, 06A07}
\keywords{Bruhat order, Weyl groups, log-concavity, Mahonian numbers, Coxeter groups}

\maketitle

\section{Introduction}
\label{sec:intro}

Let $\W$ be a finite Weyl group with simple reflections $S$, length function
$\ell$, and Bruhat order $\le$. For $u\le v$ in $\W$, the \emph{Bruhat
interval} $[u,v]=\{z\in \W : u\le z\le v\}$ carries a natural graded rank
sequence
\[
a_k([u,v]) \;=\; \#\{z\in[u,v] : \ell(z)=\ell(u)+k\}, \qquad 0\le k\le
\ell(u,v):=\ell(v)-\ell(u),
\]
and $[u,v]$ is \emph{rank log-concave} if $a_k^2\ge a_{k-1}a_{k+1}$ for every
$1\le k\le \ell(u,v)-1$. Brenti asked, as Conjecture 2.11 of his 2024 survey
of open problems on Coxeter groups and unimodality~\cite{brenti2024open},
whether this holds for \emph{every} Bruhat interval of \emph{every} finite
Weyl group. The question is delicate: the log-concavity statement is false
for general finite Coxeter groups, failing already for the non-crystallographic
group $H_3$ (Section~\ref{sec:prelim}), so any positive answer for Weyl
groups specifically must use the crystallographic structure in an essential
way, not merely finiteness of the group.

At the time of writing, the conjecture was verified only for a modest list of
small-rank cases: types $A_n$ and $D_n$ for $n\le 5$, $B_n$ for $n\le 4$,
$B_5$ restricted to intervals of length at least $20$, $F_4$, and the
dihedral groups~\cite{brenti2024open}. Beyond direct verification, the
strongest structural result in the literature is the top-heaviness inequality
of Bj\"orner and Ekedahl~\cite{bjorner-ekedahl} for parabolic lower intervals,
which constrains the \emph{shape} of the rank sequence but does not compare
ratios $a_k^2/(a_{k-1}a_{k+1})$ across different intervals and so neither
implies nor suggests any statement about \emph{where} the log-concavity ratio
is smallest.

\subsection*{Contributions} We report four results.

\begin{enumerate}[label=(\roman*)]
\item An exhaustive computational verification that pushes the confirmed
frontier well past Brenti's list: every Bruhat interval is checked, in exact
integer arithmetic, in types $A_2$--$A_7$, $B_2$--$B_6$ (the \emph{full}
range, not only $\ell(u,v)\ge 20$), $D_4$--$D_6$, $E_6$, $F_4$, and $G_2$ —
$1{,}079{,}490{,}991$ intervals total, zero violations
(Theorem~\ref{thm:V1}). We complement this with non-exhaustive evidence at
higher rank: near-top slab sweeps reaching $A_{10}$, $D_8$, and $E_7$
(apparently the first log-concavity data of any kind recorded for $E_7$), and
$124{,}944$ sampled and targeted intervals in $B_7$, $B_8$, $D_7$, $E_7$.

\item An extremal conjecture, apparently new: in every irreducible,
simply-laced Weyl group, the global minimum of the log-concavity ratio over
\emph{all} Bruhat intervals is attained by the single longest interval
$[e,w_0]$ (Conjecture~\ref{conj:F1}). This is consistent with every exhaustive
and near-top data point we have, including two structural counterexamples we
found along the way showing that both hypotheses --- irreducibility and the
simply-laced restriction --- are necessary.

\item In type $A$, a determination --- conditional on a single explicitly
stated open lemma --- of how the minimum log-concavity ratio of the
Mahonian distribution (the $[e,w_0]$ interval of $S_m$) approaches $1$,
extending the local central-limit-theorem technique of Canfield, Janson,
and Zeilberger~\cite{cjz2011} from the central coefficient of the Gaussian
binomial to the global minimum and an explicit constant. A local limit
theorem pinning the ratio near the center is unconditionally proved for
every $m$; the minimum itself is determined exactly, in integer
arithmetic, for every $4\le m\le560$; and the sharp asymptotic
$\sigma_m^2(r_m-1)=1-\tfrac{27}{25}m^{-1}+O(m^{-2})$ is proved conditional
on one explicitly stated lemma about the exponentially tilted Mahonian
distribution (Conjecture~\ref{conj:CL}), needed only for $m\ge561$, with
the reduction complete and every constant explicit; the lemma is in turn
reduced to four explicitly stated open cumulant statements
(Section~\ref{sec:F2}).

\item A classification of the observed equality cases --- log-concavity
holding with equality --- as rank-two dihedral parabolic patterns of order
$m\in\{4,6\}$, together with a candidate explanation, via the
crystallographic restriction theorem, for why Weyl groups escape the $H_3$
failure: the $H_3$ counterexample sits on a dihedral core of order $m=5$,
which is precisely the order excluded from crystallographic reflection
groups (Section~\ref{sec:F3}).
\end{enumerate}

\subsection*{Related work} Beyond Brenti's survey and Bj\"orner--Ekedahl,
several recent papers touch adjacent territory without overlapping our
results. Burrull, Gui, and Hu~\cite{burrull-gui-hu} prove an asymptotic
log-concavity result for dominant lower intervals in \emph{affine} Weyl
groups under dilation, a genuinely different regime (affine, parabolic,
asymptotic-in-dilation) from the finite, full-interval statement studied
here, and record a non-unimodal example due to Stanton showing the naive
parabolic analogue of Conjecture 2.11 can fail. Kessouri, Ahmia, Arslan, and
Mesbahi~\cite{kessouri2024} prove ordinary (finite-$n$) log-concavity of
type-$B$ $q$-Mahonian numbers, with no asymptotic statement, and so do not
overlap with our type-$A$ asymptotic (iii). We also note, as a caution about
the reliability of open-problem surveys in general rather than as a
conflict, that another conjecture from the same survey was recently refuted
by Chapelier-Laget and Fromentin~\cite{chapelier-fromentin}. A fresh
literature sweep at the time of writing (2026-08-06) found no paper stating,
proving, or refuting any of (i)--(iv).

Every claim in this note is accompanied by a reproducible verification
artifact: source code, exact-arithmetic result logs, and (where a claim is
fully proved) machine-checked numeric certificates are available at
[repository URL to be added on submission].

\section{Preliminaries}
\label{sec:prelim}

Let $\W$ be a finite Coxeter group with simple generators $S$, length
function $\ell$, and Bruhat order $\le$: $u\le v$ iff some (equivalently,
every) reduced word for $v$ contains a reduced word for $u$ as a subword.
For a finite \emph{Weyl} group $\W$, i.e.\ the group generated by
reflections in the root system of a Cartan matrix, $\W$ is realized
concretely as a subgroup of orthogonal transformations of a Euclidean space
and every construction below is computed directly from the Cartan matrix
(Section~\ref{sec:methods}).

For $u\le v$ write $[u,v]=\{z\in \W : u\le z\le v\}$ for the Bruhat interval,
graded by $\ell(z)-\ell(u)$, with rank sequence
\[
a_k([u,v]) = \#\{z\in [u,v] : \ell(z)=\ell(u)+k\}, \qquad 0\le k \le
\ell(u,v):=\ell(v)-\ell(u).
\]
All $a_k([u,v])>0$ for Bruhat intervals, so we may define the
\emph{log-concavity ratio}
\[
\rho([u,v]) = \min_{1\le k\le \ell(u,v)-1} \frac{a_k([u,v])^2}{a_{k-1}([u,v])\,
a_{k+1}([u,v])} \qquad (\ell(u,v)\ge 2),
\]
and $[u,v]$ is rank log-concave exactly when $\rho([u,v])\ge 1$. Brenti's
Conjecture 2.11~\cite{brenti2024open} asserts $\rho([u,v])\ge 1$ for every
Bruhat interval of every finite \emph{Weyl} group (crucially: not every
finite Coxeter group, by Example~\ref{ex:H3} below).

The interval $[e,w_0]$, where $w_0$ is the longest element, is distinguished:
its rank sequence is exactly the coefficient sequence of the Poincar\'e
polynomial $W(q) = \sum_{w\in \W} q^{\ell(w)} = \prod_{i} [d_i]_q$ (product
over the degrees $d_i$ of $\W$), so $a_k([e,w_0]) = [q^k]\,W(q)$.

\begin{example}[Type $A$: Mahonian numbers]
For $\W=S_m$ (type $A_{m-1}$), $\ell$ is the number of inversions,
$W(q)=[m]_q! = \prod_{i=1}^m (1+q+\cdots+q^{i-1})$, and $a_k([e,w_0]) =
I_m(k) := \#\{\sigma\in S_m : \inv(\sigma)=k\}$, the classical Mahonian
numbers. Log-concavity of $(I_m(k))_k$ itself is classical
(\cite{bona2004}; also a consequence of~\cite{hoggar1974,kook2006} applied to
the factorization of $[m]_q!$); our contribution in Section~\ref{sec:F2} is
the asymptotic behaviour of the \emph{minimum ratio}, not log-concavity
itself.
\end{example}

\begin{example}[Failure for general Coxeter groups: $H_3$]
\label{ex:H3}
We have independently confirmed the following directly against the primary
source. Let $(W,S)$ be of type $H_3$, $S=\{s_1,s_2,s_3\}$ with
$m(s_1,s_2)=5$ and $m(s_2,s_3)=3$, and let $u=s_3$,
$v=s_1s_2s_3s_2s_1s_2s_1s_3$. Then~\cite[Conjecture~2.11]{brenti2024open}
records the rank generating function of $[u,v]$ as
\[
1+3t+5t^2+7t^3+10t^4+10t^5+5t^6+t^7,
\]
i.e.\ rank sequence $(a_0,\dots,a_7)=(1,3,5,7,10,10,5,1)$. This sequence is
unimodal but not log-concave: at $k=3$, $a_3^2=49 < a_2a_4=50$, a margin of
$-1$. We verified this quotation directly against the text of
\cite{brenti2024open} (\emph{Some open problems on Coxeter groups and
unimodality}, to appear in Proc.\ Sympos.\ Pure Math.\ \textbf{110}; page
$10$ of the arXiv:2410.09897 PDF, Conjecture~2.11), where it appears
verbatim, including the stated Coxeter parameters
$m(s_1,s_2)=5,\,m(s_2,s_3)=3$; the arithmetic $7^2=49<5\cdot10=50$ checks
out exactly. The rank-2 parabolic
$\langle s_1,s_2\rangle$ generating the failure is dihedral of order
$2m(s_1,s_2)=10$, i.e.\ an $I_2(5)$ core — this is the fact exploited in
Section~\ref{sec:F3}.
\end{example}

\section{Computational methodology}
\label{sec:methods}

All computations use exact integer (arbitrary-precision) arithmetic; no
floating point is involved in any claim labeled a theorem. Three
independent implementations were used, cross-checked against each other and
against known invariants before any new result was trusted:

\begin{itemize}
\item A generic Cartan-matrix engine (\texttt{weyl.py}), which builds any
finite Weyl group directly from its Cartan matrix and computes $[u,v]$ by
bitset-encoded up/down Bruhat-order BFS; used for the exhaustive tier.
\item A complement-BFS engine for near-top intervals (\texttt{scaled.py},
\texttt{scaled\_general.py}), which computes rank sequences of lower
intervals $[e,v]$ with $v$ close to $w_0$ by BFS on the (small) complement
$\W\setminus [e,v]^{\downarrow}$ rather than on $[e,v]$ itself, making groups
far too large to enumerate exhaustively tractable in a thin slab below $w_0$.
\item A root-action engine (\texttt{fast.py}) for high-throughput sampling
and seeded perturbation search.
\end{itemize}

Each engine was validated against independent invariants before use: known
group order $|\W|$; length $\equiv$ number of inversions (root-system
definition) for every enumerated element; Poincar\'e polynomial $\equiv$ the
degree-product formula $\prod_i[d_i]_q$; and, where two engines cover the
same group, engine-versus-engine agreement on every interval.

We report results in three tiers, and state precisely what each does and
does not establish:

\begin{description}
\item[Exhaustive] Every Bruhat interval of the group is enumerated and
checked. This is the only tier that constitutes a full verification of
Conjecture 2.11 for the group in question (Theorem~\ref{thm:V1},
Table~\ref{tab:exhaustive}).
\item[Near-top slab] Only lower intervals $[e,v]$ with $\ell(w_0)-\ell(v)$
below a fixed bound are checked; this is evidence for the conjecture and for
Conjecture~\ref{conj:F1} in the checked groups, not a full verification.
\item[Sampling / seeded] Random or targeted samples; coverage evidence only.
\end{description}

\begin{remark}[$E_6$/$B_6$ segmentation and a lost witness]
\label{rem:e6segment}
The $E_6$ and $B_6$ exhaustive runs were each executed as two complementary,
non-overlapping segments (partitioned by the bottom element $u$) after an
earlier single-pass attempt was interrupted by a process kill; the segment
boundaries are documented and their union is a genuine full cover of the
group (verified: the segment interval counts sum exactly to the totals in
Table~\ref{tab:exhaustive}). For $B_6$ this recovery is additionally
corroborated by an independent non-segmented full run, so no data is in
question. For $E_6$, the first segment ($u<6000$) survived only as a
running-minimum checkpoint stream, without the interval $[u,v]$ or rank
sequence attaining that minimum; we therefore report the exact minimum ratio
for $E_6$ but cannot exhibit its witness interval in this note, and flag a
re-scan of $u<6000$ as a documented, low-cost follow-up rather than a gap
that affects the truth of Theorem~\ref{thm:V1}.
\end{remark}

All raw result logs, per-run dossiers, and the CI configuration used to
reproduce the exhaustive tier are included in the accompanying repository
([repository URL to be added on submission]).

\section{Verification results}
\label{sec:results}

\begin{table}[htbp]
\centering
\caption{Exhaustive tier: every Bruhat interval, exact integer arithmetic.}
\label{tab:exhaustive}
\small
\begin{tabular}{@{}lrrrl@{}}
\toprule
Group & $|W|$ & \#intervals ($\ell\ge2$) & min ratio & witness (index $k$) \\
\midrule
$A_2$ & 6 & 19 & 2.000000 & $[e,121]$, $k=1$ \\
$A_3$ & 24 & 213 & 1.388889 & $[e,12321]$, $k=2$ \\
$A_4$ & 120 & 3{,}781 & 1.210000 & $[e,w_0]$, $k=5$ \\
$A_5$ & 720 & 98{,}407 & 1.122222 & proper (ties $[e,w_0]$), $k=7$ \\
$A_6$ & 5{,}040 & 3{,}550{,}919 & 1.079096 & proper (ties $[e,w_0]$), $k=10$ \\
$A_7$ & 40{,}320 & 170{,}288{,}585 & 1.054250 $=919681/872356$ & proper, $\ell=27$ (ties $[e,w_0]$), $k=14$ \\
$B_2$ & 8 & 33 & 1.000000 & $[e,1212]$, $k=2$ \\
$B_3$ & 48 & 847 & 1.000000 & $[e,2323]$, $k=2$ \\
$B_4$ & 384 & 40{,}249 & 1.000000 & $[e,3434]$, $k=2$ \\
$B_5$ & 3{,}840 & 3{,}089{,}459 & 1.000000 & $[e,4545]$, $k=2$ \\
$B_6$ & 46{,}080 & 350{,}676{,}009 & 1.000000 & $[e,5656]$, $k=2$ \\
$D_4$ & 192 & 9{,}817 & 1.136232 & proper (ties $[e,w_0]$), $\ell=11$, $k=5$ \\
$D_5$ & 1{,}920 & 745{,}377 & 1.069459 & proper (ties $[e,w_0]$), $\ell=19$, $k=10$ \\
$D_6$ & 23{,}040 & 84{,}339{,}681 & 1.040703 & proper, $\ell=27$ (ties $[e,w_0]$), $k=15$ \\
$E_6$ & 51{,}840 & 466{,}250{,}713 & 1.028446 & witness unrecoverable, see Remark~\ref{rem:e6segment} \\
$F_4$ & 1{,}152 & 396{,}809 & 1.000000 & $[e,2323]$, $k=2$ \\
$G_2$ & 12 & 73 & 1.000000 & $[e,1212]$, $k=2$ \\
\bottomrule
\end{tabular}
\end{table}

\begin{theorem}[Exhaustive tier]
\label{thm:V1}
Let $\W$ be one of the Weyl groups $A_2,\dots,A_7$, $B_2,\dots,B_6$,
$D_4,D_5,D_6$, $E_6$, $F_4$, $G_2$. Then every Bruhat interval
$[u,v]\subseteq \W$ is rank log-concave. The verification is exhaustive:
every interval of length $\ge 2$ was enumerated and checked in exact
integer arithmetic (Table~\ref{tab:exhaustive}), totalling
$1{,}079{,}490{,}991$ intervals with zero violations.
\end{theorem}

\begin{remark}
Relative to the previously recorded frontier (Brenti's list: $A_n,D_n$ for
$n\le 5$; $B_n$ for $n\le 4$; $B_5$ only for $\ell(u,v)\ge 20$; $F_4$;
dihedral groups), the cases $A_6$, $A_7$, $B_5$ (full range), $B_6$, $D_6$,
$E_6$ are new.
\end{remark}

\begin{table}[htbp]
\centering
\caption{Near-top slab tier: lower intervals $[e,v]$, $\ell(w_0)-\ell(v)\le
c$ (not exhaustive). Exact fractions for the min ratio, where computed,
are given in the text.}
\label{tab:neartop}
\footnotesize
\begin{tabular}{@{}lrlrrl@{}}
\toprule
Group & $|\W|$ & slab & \#checked & min ratio & at $k$ \\
\midrule
$A_7$ & 40{,}320 & $c\le4$ & full & 1.054250 & $[e,w_0]$, $14$ \\
$A_8$ & 362{,}880 & $c\le4$ & full & 1.038942 & $[e,w_0]$, $18$ \\
$A_9$ & 3{,}628{,}800 & $c\le3$ & 209/209 & 1.028950 & $[e,w_0]$, $22$ \\
$A_{10}$ & 39{,}916{,}800 & $c\le2$ & 11/65\textsuperscript{*} & 1.022102 & tie\textsuperscript{\dag}, $27$ \\
$D_7$ & 322{,}560 & $c\le3$ & 112/112 & 1.025574 & $[e,w_0]$, $21$ \\
$D_8$ & 5{,}160{,}960 & $c\le3$ & 156/156 & 1.017122 & $[e,w_0]$, $28$ \\
$E_7$ & 2{,}903{,}040 & $c\le2$ & full & 1.011829 & $[e,w_0]$, $31$ \\
\bottomrule
\end{tabular}
\\[2pt]
\footnotesize\textsuperscript{*}partial, resumable; \textsuperscript{\dag}$[e,w_0]$ attains
the minimum and one proper interval exactly ties it.
\end{table}

\begin{theorem}[Near-top tier --- lower intervals only; not exhaustive]
\label{thm:V2}
For each pair $(\W,c)$ below, every lower interval $[e,v]$ with
$\ell(w_0)-\ell(v)\le c$ is rank log-concave; moreover the minimum of
$\rho([e,v])$ over this slab is attained at $v=w_0$
(Table~\ref{tab:neartop}):
\[
(A_7,4),\ (A_8,4),\ (A_9,3),\ (D_7,3),\ (D_8,3),\ (E_7,2).
\]
In $A_{10}$, $[e,w_0]$ and a partial slab ($11$ of $65$ candidates at
$c=2$) were checked, all log-concave, with one proper interval exactly
tying $\rho([e,w_0])$.
\end{theorem}

\begin{remark}
The two exact rational minima recorded in Table~\ref{tab:neartop}:
$\rho([e,w_0])=919681/872356$ in $A_7$, and $\rho([e,w_0])=65523/64757$ in
$E_7$; the latter's rank sequence sums to exactly $2{,}903{,}040=|W(E_7)|$,
confirming it is the full interval and not a shorter proper one.
\end{remark}

\begin{remark}
Theorem~\ref{thm:V2} checks only lower intervals in a thin slab below
$w_0$; it is evidence for Conjecture 2.11 and for
Conjecture~\ref{conj:F1} in these groups, not a verification of either
beyond the stated slab. $E_7$ is, to our knowledge, the first
log-concavity verification data of any kind recorded for that group.
\end{remark}

\begin{table}[htbp]
\centering
\caption{Sampling / seeded tier (not exhaustive; coverage counts).}
\label{tab:sampling}
\small
\begin{tabular}{@{}llrlll@{}}
\toprule
Probe & Group & \#intervals & design & min ratio & note \\
\midrule
random (seed 7) & $B_7$ & 20{,}000 & $4\le\ell(u,v)\le12$ & 1.000000 & $(1,2,2,2,1)$ pattern \\
random (seed 7) & $D_7$ & 20{,}000 & same & $1.157143=81/70$ & proper, $k{=}4$ \\
random (seed 7) & $E_7$ & 20{,}000 & same & $1.142400=14641/12816$ & proper, $k{=}5$ \\
seeded (seed 3) & $B_7$ & 1{,}324 & $m{=}4$ core $+$ 0--6 covers & 1.000000 only at pure core & 28 ratio-1 hits \\
seeded (seed 4) & $B_7$ & 32{,}458 & $m{=}4$ core $+$ 0--10 covers & pert.$\ge1$: $\ge 1.079153$ & margins 4--8 \\
seeded (seed 4) & $B_8$ & 31{,}162 & $m{=}4$ core $+$ 0--10 covers & pert.$\ge1$: $\ge 1.087990$ & margins 4--8 \\
\bottomrule
\end{tabular}
\end{table}

\begin{remark}
The seeded probe attempts a fixed number of perturbations of the $m=4$
dihedral core (the header count of each seeded run above: $4{,}000$;
$100{,}000$; $100{,}000$); an attempt is discarded, uncounted, whenever the
targeted perturbation does not realize a genuine Bruhat up-cover chain. The
counts reported in the \#intervals column and used throughout this paper
are the surviving, actually-checked intervals only.
\end{remark}

\begin{proposition}[Sampling tier --- not exhaustive; coverage statement only]
\label{thm:V3}
\textup{(i)} $60{,}000$ Bruhat intervals ($20{,}000$ each in $B_7$, $D_7$,
$E_7$), sampled uniformly from a random-walk ensemble with
$4\le\ell(u,v)\le 12$ (seed $7$, reproducible), are all rank log-concave.
\textup{(ii)} $64{,}944$ seeded intervals containing an $m=4$ dihedral
braid core perturbed by $0$--$10$ extra cover steps are all rank
log-concave; every perturbation that destroys the pure dihedral pattern
strictly raises the ratio (Table~\ref{tab:sampling}; the seeded search
attempted more perturbations than this, discarding any that did not
realize a genuine cover chain --- see the remark following the table).
\end{proposition}

In total, the sampling and seeded probes cover $124{,}944$ intervals with
zero violations and zero equality cases outside the $(1,2,\dots,2,1)$
pattern classified in Section~\ref{sec:F3}.

\section{An extremal conjecture}
\label{sec:F1}

\begin{conjecture}[F1]
\label{conj:F1}
Let $\W$ be a finite Weyl group that is \emph{irreducible} and
\emph{simply-laced} (type $A_n$, $D_n$, $E_6$, $E_7$, or $E_8$), with
longest element $w_0$. Then
\[
\min_{\substack{[u,v]\subseteq \W \\ \ell(u,v)\ge 2}} \rho([u,v])
\;=\; \rho([e,w_0])
\;=\; \min_{1\le k\le \ell(w_0)-1} \frac{b_k^2}{b_{k-1}b_{k+1}},
\]
where $b_k=[q^k]\,W(q)$: the global minimum of the log-concavity ratio over
\emph{all} Bruhat intervals of $\W$ is attained by the full interval
$[e,w_0]$, at a central index. Proper intervals may tie this minimum
(observed for $A_5,A_6,A_7,D_6$ in the exhaustive tier, and for a further
proper interval in the $A_{10}$ near-top slab) but never fall below it.
\end{conjecture}

Both hypotheses in Conjecture~\ref{conj:F1} are necessary; we record the two
sharp counterexamples that force this wording.

\begin{example}[Simply-laced is necessary]
\label{ex:F1-simplylaced}
In $B_3$, the interval $[e,(s_2s_3)^2]$ has rank sequence $(1,2,2,2,1)$ and
$\rho=2^2/(2\cdot2)=1$ exactly, on a plateau, strictly below
$\rho([e,w_0(B_3)])=8/7$. The analogous $B_4$ interval
$[e,(s_3s_4)^2]$ gives $\rho=1<968/897=\rho([e,w_0(B_4)])$. Note
$[e,(s_2s_3)^2]$ is \emph{rationally smooth} (its rank sequence is
palindromic), so dropping the simply-laced hypothesis falsifies the
conjecture even within the rationally-smooth subclass discussed below.
\end{example}

\begin{example}[Irreducible is necessary --- apparently new]
\label{ex:F1-irreducible}
Let $\W=A_1\times D_4$ (simply-laced, reducible) and $v=(e,w_0(D_4))\in\W$.
The interval $[e,v]$ is isomorphic, as a graded poset, to $[e,w_0(D_4)]$
inside the $D_4$ factor, with rank sequence $P_{D_4}=
(1,4,9,16,23,28,30,28,23,16,9,4,1)$ and
$\rho([e,v])=28^2/(23\cdot30)=392/345=1.1362318\ldots$ at $k=5$. This is a
\emph{proper} interval of $\W$: $\ell(v)=12<13=\ell(w_0(\W))$. The full
interval has rank sequence $[2]_q\cdot P_{D_4}=
(1,5,13,25,39,51,58,58,51,39,25,13,5,1)$ and
$\rho([e,w_0(\W)])=58^2/(51\cdot58)=58/51=1.1372549\ldots$ Since
$392\cdot51=19992<20010=345\cdot58$,
\[
\rho([e,v]) \;=\; \frac{392}{345} \;<\; \frac{58}{51} \;=\; \rho([e,w_0(\W)]),
\qquad \text{gap } \frac{58}{51}-\frac{392}{345}=\frac{2}{1955}\approx0.00102,
\]
so a \emph{proper} Bruhat interval of the reducible simply-laced group
$A_1\times D_4$ strictly beats $[e,w_0]$: the extremal principle fails
outright once reducibility is allowed, not merely on some restricted
subclass. (The witness $v$ happens to be smooth, but this already refutes
the unrestricted statement for reducible $\W$, independently of
smoothness.) An exhaustive check of all $27$ reducible simply-laced types of
rank $\le 6$ found $A_1\times D_4$ to be the \emph{only} violator.
\end{example}

\begin{remark}[Status: conjectural, not proved]
Conjecture~\ref{conj:F1} itself --- for arbitrary $[u,v]$, not merely
rationally smooth ones --- is not proved for any infinite family. What
supports it is: exhaustive verification with zero violations for every
irreducible simply-laced group in the exhaustive tier
($A_2$--$A_7$, $D_4$--$D_6$, $E_6$; Table~\ref{tab:exhaustive}), in every
case with $[e,w_0]$ attaining (and, in $A_5,A_6,A_7,D_6$, exactly tied by a
specific proper interval); and near-top slab confirmation with no exception
in $A_8$--$A_{10}$, $D_7$--$D_8$, $E_7$ (Table~\ref{tab:neartop}). Neither
tier constitutes a proof for general rank.
\end{remark}

\begin{theorem}[F1 restricted to rationally smooth intervals]
\label{thm:F1-smooth}
Let $\W$ be a finite irreducible simply-laced Weyl group of rank $\le 6$
($A_1,\dots,A_6,D_4,D_5,D_6,E_6$). For every rationally smooth $v\in\W$
--- i.e.\ every $v$ for which the Poincar\'e polynomial $P_{[e,v]}(q)$
factors as a product of $q$-integers~\cite{carrell1994} --- $\rho([e,v])\ge
\rho([e,w_0])$, with equality iff $v=w_0$. Unconditionally, this also holds
for every $m$ in type $A_{m-1}$ with $m\le17$, modulo the classical
rational-smoothness factorization fact cited above, which our own working
notes verified computationally only through $m\le7$ rather than re-deriving
in full generality; we flag this explicitly rather than presenting the
$m\le17$ range as resting on nothing beyond our own work.
\end{theorem}

\begin{conjecture}[Staircase domination]
\label{conj:staircase}
In type $A_{m-1}$, the exponent multiset of the rationally-smooth
Poincar\'e polynomial $P_v$ is always a gap-free (``staircase'') multiset
dominated (coefficientwise, after sorting) by the degree multiset
$\{2,\dots,m\}$.
\end{conjecture}

\begin{remark}
Conjecture~\ref{conj:staircase} is verified without exception on
$91{,}355$ dominated pairs, and, in the specific instances
Theorem~\ref{thm:F1-smooth} needs, on a further $131{,}036$ pairs for
$m\le17$; it has not been proved for general $m$. Granting it, the
$m\le17$ restriction in Theorem~\ref{thm:F1-smooth}'s type-$A$ clause can
be removed, giving the same inequality for every $m$.
\end{remark}

\begin{remark}
Theorem~\ref{thm:F1-smooth} is a genuine partial proof, but of a
\emph{strictly weaker} statement than Conjecture~\ref{conj:F1}: it compares
$\rho([e,v])$ against $\rho([e,w_0])$ only for $v$ rationally smooth, not
for arbitrary $v$. Both Examples~\ref{ex:F1-simplylaced} and
\ref{ex:F1-irreducible} witness this restricted statement's own sharpness
(the $B_3$ counterexample is itself rationally smooth), so the same
irreducible-simply-laced hypotheses are necessary here too. Type $D_n$ for
$n\ge7$ and types $E_7,E_8$ remain open even for this restricted subclass.
This theorem is the paper's best current route to a proved instance of the
extremal principle, and we record it as the natural target for future work
on Conjecture~\ref{conj:F1} in full.
\end{remark}

\section{The type-\texorpdfstring{$A$}{A} case: Mahonian asymptotics}
\label{sec:F2}

Throughout this section $\W=S_m$ (type $A_{m-1}$), $I_m(k)=\#\{\sigma\in
S_m:\inv(\sigma)=k\}$ so that $a_k([e,w_0])=I_m(k)$, $N=\binom m2$,
$\sigma^2=\Var(\inv)=m(m-1)(2m+5)/72$, $r_m(k)=I_m(k)^2/(I_m(k-1)I_m(k+1))$,
and $r_m=\min_{1\le k\le N-1}r_m(k)$. Write $y=(k-N/2)/\sigma$, and let
$B_m:=(S_4-m)/(240\sigma^4)$ where $S_r=\sum_{j=1}^m j^r$, so $B_m=\tfrac{27}{25}m^{-1}(1+O(m^{-1}))$.

\subsection*{What is proved unconditionally}
The following local limit theorem for the log-concavity ratio, in the
classical style of~\cite{petrov1975}, pinning it
near the center for \emph{every} $m$, is fully proved --- with every
constant explicit and both an analytic argument and an exact finite check
covering every case, and independently referee-audited (mathematics and
numerics separately, in two independent passes), with the referees'
verdict ``survives with minor repairs'' and those repairs applied.

\begin{theorem}[Pointwise and window law for the Mahonian ratio]
\label{thm:G1}
Let $p(k)=I_m(k)/m!$, $Z(y)=(2\pi\sigma^2)^{-1/2}e^{-y^2/2}$, and let
$\mathrm{He}_4$ denote the fourth probabilists' Hermite polynomial.
\begin{enumerate}[label=(\alph*)]
\item \textup{(Pointwise.)} For every $m\ge4$ and every $0\le k\le N$,
\[
\Bigl| p(k) - Z(y)\bigl[1-\tfrac{B_m}{12}\mathrm{He}_4(y)\bigr] \Bigr|
\;\le\; \frac{0.45}{\sigma m^2}.
\]
\item \textup{(Window ratio law.)} For every fixed $y_0\ge0$ there are
explicit $C_2(y_0)$ and $m_1(y_0)$ (e.g.\ $C_2(1)=3.1$, $m_1(1)=180$; the
constant grows with $y_0$, e.g.\ $C_2(3)=3940$) such that for every $m\ge
m_1(y_0)$ and every $k$ with $|y|\le y_0$,
\[
\sigma^2\log r_m(k) \;=\; 1+B_m(y^2-1)+E_1(k), \qquad |E_1(k)|\le
\frac{C_2(y_0)}{m^2}.
\]
Moreover (a) holds, with the same displayed constant, for every $4\le
m\le109$ by direct exact-arithmetic verification, and the window law of
(b) holds, with the same displayed constants, for every $4\le m\le229$;
both are also fully covered by the analytic argument above outside these
ranges, so neither part carries a threshold in $m$ --- the exact-arithmetic
ranges are an independent check, not the limit of what is proved, and the
window law in (b) is unconditional (no asymptotic-only asterisk) throughout
the range any downstream use in this paper requires.
\end{enumerate}
\end{theorem}

Theorem~\ref{thm:G1} determines the central ratio exactly to the stated
order: at $y=0$, part (b) gives
\begin{equation}
\label{eq:central-ratio}
\sigma^2\bigl(r_m(\lfloor N/2\rfloor)-1\bigr) \;=\; 1-\tfrac{27}{25}m^{-1}+O(m^{-2}),
\end{equation}
which we have independently confirmed to four significant digits against
exact Mahonian data by $m=50$.

The second unconditional ingredient is exact: the global minimum itself
--- its location, its value, and the sharp finite-$m$ bound --- has been
computed exactly for every $4\le m\le560$.

\begin{theorem}[Exact finite range]
\label{thm:finite560}
For every $5\le m\le560$, with every verdict computed in exact integer or
rational arithmetic:
\begin{enumerate}[label=(\roman*)]
\item the minimum of $r_m(k)$ over $1\le k\le N-1$ is attained at the
central index $k=\lfloor N/2\rfloor$ (for $N$ odd, the mirror index
$\lceil N/2\rceil$ ties it exactly, as palindromy forces);
\item consequently $r_m=r_m(\lfloor N/2\rfloor)$;
\item $\sigma^2(r_m-1)\ge\tfrac{187}{216}$, with equality if and only if
$m=6$;
\item $m\mapsto\sigma^2(r_m-1)$ is strictly increasing for $6\le m\le560$,
reaching $\sigma^2(r_m-1)=0.998072\ldots$ at $m=560$, where
$m\bigl(1-\sigma^2(r_m-1)\bigr)\approx1.07935$, consistent with the limit
$\tfrac{27}{25}=1.08$.
\end{enumerate}
At $m=4$, and only there in the computed range, the minimizing index is
not central ($\mathrm{argmin}=2$, $N/2=3$) and
$\sigma^2(r_4-1)=\tfrac{91}{108}<\tfrac{187}{216}$.
\end{theorem}

The computation is a checkpointed harness over the $q$-factorial
recurrence for $I_m$; every verdict is an integer or exact-rational
comparison, floating point appears only in printed displays, and the
complete result logs --- one row for every $4\le m\le560$, all passing ---
are archived in the repository.

\subsection*{From the center to the global minimum: the reduction}
Equation~\eqref{eq:central-ratio} pins the ratio at and near the center,
and Theorem~\ref{thm:finite560} settles every $4\le m\le560$ outright. What
remains for the sharp asymptotic is to rule out the ratio dipping below
its central value somewhere in the bulk for $m\ge561$. This has now been
carried out in full, with every constant explicit, \emph{modulo exactly
one statement}: an estimate for the exponentially tilted Mahonian
distribution in the deep-tilt regime, stated precisely as
Conjecture~\ref{conj:CL} below. We first fix the frame.

For $\lambda\in\mathbb R$ let $U^\lambda_1,\dots,U^\lambda_m$ be
independent, with $U^\lambda_j$ supported on $\{0,1,\dots,j-1\}$ and
$\Pr(U^\lambda_j=i)\propto e^{-\lambda i}$, and set
$S_\lambda:=\sum_{j=1}^m U^\lambda_j$ (not to be confused with the power
sums $S_r$ above), $\mu(\lambda):=\mathbb E S_\lambda$.
At $\lambda=0$ this is the inversion statistic:
$\Pr(S_0=k)=I_m(k)/m!$. Since $\mu'(\lambda)=-\Var S_\lambda<0$, every
interior $k$ ($1\le k\le N-1$) has a unique \emph{mean-matching tilt}
$\lambda(k)$ with $\mu(\lambda(k))=k$; write
$s^2=s^2(k):=\Var S_{\lambda(k)}$, let $\kappa_3,\kappa_4$ denote the
third and fourth cumulants of $S_{\lambda(k)}$, and let
$\varphi(t):=\mathbb E\,e^{it(S_{\lambda(k)}-k)}$. By palindromy
$r_m(N-k)=r_m(k)$ and $\lambda(N-k)=-\lambda(k)$, so we take $k\le N/2$,
i.e.\ $\lambda(k)\ge0$, throughout, and set $\omega:=m\lambda(k)$.

\begin{conjecture}[Core lemma $\mathrm{CL}(79,20,0.89)$]
\label{conj:CL}
For every $m\ge561$ and every interior $k$ with $s^2(k)\ge79$ and
$4/m<\lambda(k)\le0.89$,
\[
r_m(k)-1 \;=\; \frac{1}{s^2}\Bigl(1+\theta\,\frac{20}{\min(m,s^2)}\Bigr)
\qquad\text{for some }\theta=\theta(m,k)\in[-1,1].
\]
\end{conjecture}

\begin{remark}
Three notes on the statement. (a) The hypothesis $s^2\ge79$ is part of
the frozen specification of the reduction but never binds: on the stated
tilt range $s^2\ge1122800/7921>141$ is proved. (b) Only the lower-bound
half ($\theta\ge-1$) is consumed by Theorem~\ref{thm:F2}. (c) The
complementary small-tilt range $\lambda(k)\le4/m$ is not part of the
conjecture because it is covered unconditionally by the proved window law
(region (d) of Remark~\ref{rem:reduction}). The reduction behind
Theorem~\ref{thm:F2} needs this statement for every $m\ge401$; the range
$401\le m\le560$ is discharged by the exact computation of
Theorem~\ref{thm:finite560}, which is why the conjecture is stated, and
is needed, only for $m\ge561$.
\end{remark}

\begin{theorem}[F2, sharp asymptotic --- conditional]
\label{thm:F2}
Assume Conjecture~\ref{conj:CL}. Then
\[
\sigma^2(r_m-1) \;=\; 1-\tfrac{27}{25}m^{-1}+O(m^{-2}),
\qquad\text{equivalently}\qquad r_m-1 \sim \frac{36}{m^3},
\]
and in particular $\sigma^2(r_m-1)\to1$. The error term is explicit and
two-sided: for every $m\ge561$,
\[
1-B_m-\frac{C_A}{m^2} \;\le\; \sigma^2(r_m-1) \;\le\; 1-B_m+\frac{1.8}{m^2},
\qquad C_A=37997.85,
\]
the upper bound holding unconditionally for every $m\ge180$.
\end{theorem}

\begin{remark}[Shape of the reduction, and where the conjecture enters]
\label{rem:reduction}
The proof is an assembly in which every constant is explicit and every
step is archived, with machine-verified certificates, in the repository;
we record its shape here. Take $m\ge561$ and (WLOG) $k\le N/2$. The
interior indices split into four regions by the size of $k$ and of
$\omega(k)=m\lambda(k)$, with $K_c:=\min(\tfrac7{10}m,\,m-1)$ for
$m<1581$ and $K_c:=m-1$ beyond:
\begin{enumerate}[label=(\alph*)]
\item $k=1$: an exact pentagonal-recurrence bound gives
$r_m(1)-1\ge\tfrac{m-1}{2(m+1)}$, hence $\sigma^2(r_m(1)-1)\ge10^5$;
\item $2\le k\le K_c$: an elementary two-term bound gives
$r_m(k)-1\ge\tfrac{m-1}{2k(m+k)}$, hence $\sigma^2(r_m(k)-1)\ge1879$;
\item $k>K_c$ with $\omega(k)>4$ (the deep-tilt region): here
$\min(m,s^2)\ge79.5$, $s^2\le0.72711\,\sigma^2$, and
$\lambda(k)\le\log(17/7)<0.89$, all proved; Conjecture~\ref{conj:CL} ---
consumed here and nowhere else --- gives
$\sigma^2(r_m(k)-1)\ge\bigl(1-\tfrac{20}{79.5}\bigr)/0.72711\ge1.0293$;
\item $k>K_c$ with $\omega(k)\le4$: the window law of
Theorem~\ref{thm:G1}, extended unconditionally and with explicit
constants to the full small-tilt window $|\omega|\le4$, gives
$\sigma^2(r_m(k)-1)\ge1-B_m-C_A/m^2$.
\end{enumerate}
Regions (a)--(c) all exceed $1.02>1-B_m$, so the global minimum lies in
region (d); the lower bound follows, and the upper bound is
$r_m\le r_m(\lfloor N/2\rfloor)$ together with the explicit form
of~\eqref{eq:central-ratio}. The size of $C_A$ is an artifact of one
lossy triangle-inequality step (the measured analogue of the dominant
contribution is about $5$); any fixed constant suffices for the sharp
asymptotic.
\end{remark}

\subsection*{Reduction of the core lemma to four cumulant statements}
Conjecture~\ref{conj:CL} has itself been reduced, with explicit constants
throughout, to four statements about the cumulants of the tilted law.
Partition the deep-tilt values $\omega=m\lambda\in(4,\,0.89\,m]$ into
seven bands
\[
\mathcal W_1=(4,5],\ \mathcal W_2=(5,6],\ \mathcal W_3=(6,8],\
\mathcal W_4=(8,10],\ \mathcal W_5=(10,20],\ \mathcal W_6=(20,40],\
\mathcal W_7=(40,\infty),
\]
each intersected with $\{\lambda\le0.89\}$, and normalize the cumulants
on the tilt scale:
\[
\tilde\kappa_3:=\frac{|\kappa_3|\,\lambda}{s^2},\qquad
\tilde\kappa_4:=\frac{\kappa_4\,\lambda^2}{s^2}.
\]

\begin{table}[htbp]
\centering
\caption{Per-band constants in Conjecture~\ref{conj:S} (bands of
$\omega=m\lambda$). The $J_0$ entries are six-digit displays of exact
rational constants archived in the repository.}
\label{tab:CLbands}
\small
\begin{tabular}{@{}lccccccc@{}}
\toprule
band & $\mathcal W_1$ & $\mathcal W_2$ & $\mathcal W_3$ & $\mathcal W_4$ & $\mathcal W_5$ & $\mathcal W_6$ & $\mathcal W_7$ \\
$\omega$-range & $(4,5]$ & $(5,6]$ & $(6,8]$ & $(8,10]$ & $(10,20]$ & $(20,40]$ & $(40,\infty)$ \\
\midrule
$R_3$ & 1.0 & 1.2 & 1.5 & 1.7 & 2.0 & 2.1 & 2.2 \\
$R_4$ & 0.8 & 1.4 & 2.6 & 3.5 & 5.2 & 6.0 & 6.6 \\
$C_5$ & 0.05 & 0.06 & 0.08 & 0.10 & 0.15 & 0.25 & 0.80 \\
$J_0$ & 0.682942 & 1.10268 & 1.91562 & 2.53645 & 3.66793 & 4.17806 & 4.59597 \\
\bottomrule
\end{tabular}
\end{table}

\begin{conjecture}[The four open statements]
\label{conj:S}
For every $m\ge561$ and every interior $k$ with
$\lambda=\lambda(k)\in(4/m,\,0.89]$, with $\mathcal W$ the band containing
$\omega=m\lambda$ and constants $R_3,R_4,C_5,J_0$ as in
Table~\ref{tab:CLbands}:
\begin{itemize}
\item[(S1)] \emph{(banded cumulant scales)}\quad
$\tilde\kappa_3\le R_3(\mathcal W)$ and $\tilde\kappa_4\le R_4(\mathcal W)$;
\item[(S2)] \emph{(fifth-order core remainder)}\quad for
$0\le t\le\lambda/2$, with the continuous branch of $\log\varphi$
starting at $\log\varphi(0)=0$,
\[
\Bigl|\log\varphi(t)+\frac{s^2t^2}{2}+\frac{i\kappa_3t^3}{6}
-\frac{\kappa_4t^4}{24}\Bigr| \;\le\;
C_5(\mathcal W)\,\frac{s^2t^5}{\lambda^3};
\]
\item[(S3)] \emph{(joint cancellation)}\quad
$\tilde\kappa_3^{\,2}-\tilde\kappa_4/2\le J_0(\mathcal W)$;
\item[(S4)] \emph{(a-priori ratio seed)}\quad
$\bigl|s^2\bigl(r_m(k)-1\bigr)-1\bigr|\le0.89$.
\end{itemize}
\end{conjecture}

\begin{proposition}[Reduction of the core lemma]
\label{prop:CLred}
Statements (S1)--(S4) of Conjecture~\ref{conj:S} imply
Conjecture~\ref{conj:CL}; the composed chain in fact delivers the
constant $18.23$ in place of $20$ for every $m\ge561$.
\end{proposition}

The assembly behind Proposition~\ref{prop:CLred} proceeds band by band
over Table~\ref{tab:CLbands}: (S1) and (S3) price the main
Gaussian-model term, (S2) prices the core remainder, and (S4) seeds a
self-consistency step which then closes with strict contraction; every
other ingredient --- variance floors and caps, the far-region decay of
$\varphi$, the crossover and mid-range bounds --- is proved with explicit
constants. The complete chain, together with an end-to-end machine
re-verification of its composed constant (a mixture of exact rational
certificates and high-precision numerical evaluation), is archived in
the repository.

\begin{remark}[Status and evidence for (S1)--(S4)]
\label{rem:S-evidence}
All four statements are open; none has a proof, and all four are
load-bearing --- the assembly consumes each of them, and removing any one
leaves no theorem. The list has also moved in both directions under
scrutiny: two earlier hypotheses at this level were proved outright and
retired, while (S3) and (S4) entered the list when simpler plans were
found to fail. We therefore report the numerical evidence with its
margins rather than with optimism.
\begin{itemize}
\item (S1): the measured suprema of $\tilde\kappa_3$ and $\tilde\kappa_4$
at the deep corner of $\mathcal W_7$ are $2.1215$ and $6.3552$
(large-$m$ limits $2.1303$ and $6.4113$) against the caps $2.2$ and
$6.6$ --- headroom of only $3.7\%$ and $3.9\%$, the thinnest margins
anywhere in the reduction.
\item (S2): measured remainder ratios lie between $0.0083$ and $0.2104$;
the slack over $C_5$ ranges from a factor $1.6$ to a factor $8$
depending on the band (the lowest band's slack was measured at interior
probe points, not uniformly to its edge).
\item (S3): worst measured margin $32.6\%$, at $(m,\omega)=(561,5.0)$.
This statement cannot be dispensed with: there is a point consistent
with (S1) and with $\kappa_4\ge0$ at which the estimate (S3) feeds
fails by $43\%$, so a joint bound of this shape is genuinely required
--- an earlier plan that priced the two cumulants separately is
provably dead.
\item (S4): a weak a-priori form of the conclusion itself (constant
$0.89$ where Conjecture~\ref{conj:CL} concludes
$20/\min(m,s^2)\le0.036$, since $s^2\ge m$ on this range); it is listed
separately precisely because
assuming it silently would overstate what is proved. Exact verification
of the inequality of Conjecture~\ref{conj:CL} at $m=401$ and $m=402$
(below its stated range), on $260$
adversarially chosen interior indices, passes with a $17\times$ margin,
and every tilted ratio measured in this project sits far inside
$0.89$.
\end{itemize}
\end{remark}

\begin{remark}[What is claimed, and what is not]
Unconditional: Theorems~\ref{thm:G1} and~\ref{thm:finite560}, and the
upper-bound half of Theorem~\ref{thm:F2}. Conditional: the sharp
asymptotic of Theorem~\ref{thm:F2}, on Conjecture~\ref{conj:CL} and on
nothing else; and Conjecture~\ref{conj:CL} itself would follow from
(S1)--(S4) by Proposition~\ref{prop:CLred}. Open: statements (S1)--(S4).
We do not regard the margins of Remark~\ref{rem:S-evidence} --- one of
which is under $4\%$ --- as grounds for treating any of the four as close
to proved; they are open mathematics.
\end{remark}

\begin{remark}[Location and explicit constant]
Beyond the sharp order of Theorem~\ref{thm:F2}, two refinements. The
minimizing index is exactly central for every $5\le m\le560$
(Theorem~\ref{thm:finite560}); for $m\ge561$, conditional on
Conjecture~\ref{conj:CL}, the reduction localizes it to the small-tilt
window $|\lambda(k)|\le4/m$ around the center (region (d) of
Remark~\ref{rem:reduction}). Sharpening this to
$|\mathrm{argmin}-N/2|\le1$ for every $m$ is a further, separately open
combinatorial statement, equivalent to a third-difference positivity
property of $(I_m(k))_k$, verified without exception for $5\le m\le56$
but with no proof strategy identified. The explicit finite-$m$ bound
$r_m\ge 1+\tfrac{187}{216}\sigma_m^{-2}$ for all $m\ge5$ (sharp at $m=6$;
a corrected value --- the naive guess $7/8$ fails already at $m=6$) is
exact for $5\le m\le560$ (Theorem~\ref{thm:finite560}(iii)) and holds for
every $m\ge561$ conditional on Conjecture~\ref{conj:CL}, since the
explicit lower bound of Theorem~\ref{thm:F2} exceeds $\tfrac{187}{216}$
for every $m\ge537$; granting the conjecture, the sharp bound therefore
holds for every $m\ge5$, with equality only at $m=6$.
\end{remark}

\section{Equality cases}
\label{sec:F3}

\begin{observation}[F3, equality classification]
\label{obs:F3}
Let $\W$ be a finite Weyl group and $[u,v]\subseteq\W$ a Bruhat interval
with $\ell(u,v)\ge2$ in which $a_k^2=a_{k-1}a_{k+1}$ for some interior $k$.
Then, in every case observed in our exhaustive and seeded-perturbation data,
the full rank sequence of $[u,v]$ is
\[
(a_0,\dots,a_{\ell(u,v)}) \;=\; (1,2,2,\dots,2,1) \qquad (m-1 \text{ interior
2's}, \ m\ge4),
\]
i.e.\ $[u,v]$ is poset-isomorphic to the full Bruhat interval of the
dihedral group $I_2(m)$, arising from a rank-two dihedral standard
parabolic of braid order $m$. Since the crystallographic restriction
confines rank-two parabolics of Weyl groups to $m\in\{2,3,4,6\}$ (equality
requires $m\ge4$, so effectively $m\in\{4,6\}$), the realizable equality
cores are $m=4$ (type $B_2/C_2$, present in $B_n$ and $F_4$) and $m=6$
(type $G_2$); in particular no equality occurs in simply-laced types, where
rank-two parabolics have $m\le3$.
\end{observation}

We state this as an empirical observation, not a theorem or even a fully
scoped conjecture: it holds without exception in every exhaustively checked
group (Table~\ref{tab:exhaustive}; witnesses e.g.\ $B_2$ `1212', $B_3$
`2323', $B_4$ `3434', $B_5$ `4545', $F_4$ `2323', $G_2$ `1212') and in the
$64{,}944$-interval seeded perturbation probe (Table~\ref{tab:sampling}),
where the equality wall is strict: every perturbation of an $m=4$ dihedral
core by extra cover steps that breaks the pure pattern strictly raises the
ratio, with closest non-equality margins of $4$--$8\%$ in the cases
checked. We have not attempted a proof, and have not fully resolved whether
the correct scope is ``the whole rank sequence is $(1,2,\dots,2,1)$'' (the
strong form the data support wherever a full witness was recorded) or a
weaker, purely local statement about the equality index itself; a route to
a genuinely proved short-interval instance exists via the classification of
short Bruhat intervals~\cite{short-intervals}, which we have not carried
out.

\subsection*{Why Weyl groups escape the $H_3$ failure}
We offer the following as a candidate mechanism, not a theorem. The $H_3$
counterexample of Example~\ref{ex:H3} sits on an $m=5$ dihedral core: the
parabolic $\langle s_1,s_2\rangle\cong I_2(5)$ is a subsystem of the
failing interval's defining data. The crystallographic restriction excludes
exactly $m=5$ (and every $m\notin\{2,3,4,6\}$) from occurring as a rank-two
parabolic of any Weyl group. So the only equality/violation-producing
dihedral cores available inside a Weyl group are $m=4$ and $m=6$, and by
Observation~\ref{obs:F3} together with the strict perturbation wall, these
sit at ratio exactly $1$ and are never pushed below it by any perturbation
we have tested. This is consistent with, and offered as an explanation of,
why every Weyl group in our exhaustive and near-top data avoids the
$H_3$-type failure, but it is a mechanism sketch resting on an empirical
observation (Observation~\ref{obs:F3}), not a proof that no Weyl-group
interval can ever violate log-concavity.

\section{Discussion and open problems}
\label{sec:discussion}

We record what we regard as the most promising directions, roughly in order
of tractability.

\begin{enumerate}
\item \textbf{Prove Conjecture~\ref{conj:CL} --- that is, statements
(S1)--(S4) of Conjecture~\ref{conj:S}.} By Theorem~\ref{thm:F2} and
Proposition~\ref{prop:CLred}, these four cumulant statements about the
tilted Mahonian law are now the \emph{entire} distance between the sharp
asymptotic and an unconditional theorem: the far-region decay, the
crossover and mid-range bounds, the variance floors and caps, the
assembly itself, and the whole range $4\le m\le560$ are all in place with
explicit constants. The measured margins (Remark~\ref{rem:S-evidence})
locate the risk: the cumulant caps of (S1) in the deepest band have under
$4\%$ headroom, and that is where a proof attempt --- or a counterexample
search --- should begin.
\item \textbf{Extend Theorem~\ref{thm:F1-smooth} past rank 6}, and past
type $A$ past $m=17$ unconditionally by proving
Conjecture~\ref{conj:staircase}. $D_n$ for $n\ge7$ and $E_7,E_8$ are open
even for the rationally-smooth restriction.
\item \textbf{Self-dual and rank-symmetric Bruhat intervals}
(\cite{gaetzgao2020}, which classifies self-dual intervals) are a natural
place to look for further explicit equality-case or extremal examples,
complementing Observation~\ref{obs:F3}.
\item \textbf{A proved instance of Observation~\ref{obs:F3}.} The
classification of short Bruhat intervals~\cite{short-intervals} is a
plausible route to a genuinely proved (rather than empirical) equality
classification, at least for intervals below some fixed length.
\item \textbf{$E_8$.} No exhaustive or near-top data of any kind is
recorded here for $E_8$ ($|W|>696$ billion); even a thin near-top slab
would extend Conjecture~\ref{conj:F1}'s evidence to every exceptional
simply-laced type.
\item \textbf{A quantitative form of Conjecture 2.11}, e.g.\ $\rho([u,v])
\ge 1+c(\W)$ with $c(\W)$ explicit in simply-laced types decaying as
predicted by Conjecture~\ref{conj:F1} and Theorem~\ref{thm:G1}, would unify
the extremal and asymptotic statements of this paper into a single sharp
inequality.
\end{enumerate}

We also note two points of context for the reliability of this line of
work. First, log-concavity of Bruhat rank sequences is adjacent to, but
distinct from, several actively studied log-concavity and equality-case
questions --- equivariant log-concavity of flag-variety
cohomology~\cite{equivariant-logconcave}, and equality cases of related
inequalities~\cite{stanley-yan,kahn-saks} --- none of which, on inspection,
concern the specific statement of Conjecture 2.11. Second, as a caution
about treating any single open-problem survey as a stable ground truth: a
different conjecture from the same survey we cite for Conjecture 2.11 was
recently shown false~\cite{chapelier-fromentin}. A fresh literature sweep
immediately preceding submission, repeating and extending the one performed
when this project began, found nothing overlapping with any claim in this
paper (see the reproducibility repository for the full query log).

\section*{Acknowledgments}

This work was produced with substantial use of AI assistance, and we want
to be specific about its scope rather than gesture at it in passing. Large
language models (principally Claude, Anthropic, in an agentic coding
environment) wrote the great majority of the verification code, drafted
proof arguments for the results in Sections~\ref{sec:F1} and
\ref{sec:F2}, and assisted in assembling this manuscript from our working
notes. Distinct AI sessions were used adversarially against each other's
work where the project's own house rules required it: draft proofs were
checked by independent review passes (including a from-scratch numerical
audit) before we treated them as reliable, and at least three instances
where an AI-drafted claim was numerically wrong on first pass --- a sign
error, a false finite-range certificate, and a fabricated verification
result --- were caught this way and corrected, not merely trusted.

None of this substitutes for human responsibility. The authors directed
the research questions, judged which computational and proof strategies
were worth pursuing, decided what could honestly be called a theorem versus
what had to remain a conjecture, independently checked the disclosed
constants and counterexamples, and take full responsibility for every
claim in this paper, including any error that remains. We disclose the
AI involvement because we think readers are entitled to know it, not to
diffuse authorship: the choices of what to claim, what to hold back, and
what still needs proof are ours.

% Add acknowledgment of specific colleagues/advisors here if applicable.

\begin{thebibliography}{99}

\bibitem{brenti2024open}
F.~Brenti, \emph{Some open problems on Coxeter groups and unimodality},
arXiv:2410.09897.

\bibitem{bjorner-ekedahl}
A.~Bj\"orner and T.~Ekedahl, \emph{On the shape of Bruhat intervals}, Ann. of
Math. \textbf{170} (2009), 799--817. arXiv:math/0508022.

\bibitem{cjz2011}
E.~R.~Canfield, S.~Janson, and D.~Zeilberger, \emph{The Mahonian probability
distribution on words is asymptotically normal}, Adv. Appl. Math.
\textbf{46} (2011). arXiv:0908.2089.

\bibitem{bona2004}
M.~B\'ona, \emph{A combinatorial proof of the log-concavity of a famous
sequence counting permutations}, Electron. J. Combin. \textbf{11(2)} (2004),
\#N2.

\bibitem{hoggar1974}
S.~G.~Hoggar, \emph{Chromatic polynomials and logarithmic concavity}, J.
Combin. Theory Ser. B \textbf{16} (1974).

\bibitem{kook2006}
W.~Kook, \emph{On the product of log-concave polynomials} (elementary note
on product-closure of log-concavity), 2006.

\bibitem{butler1990}
L.~M.~Butler, \emph{The $q$-log-concavity of $q$-binomial coefficients}, J.
Combin. Theory Ser. A \textbf{54} (1990), 54--63.

\bibitem{sagan1992}
B.~E.~Sagan, \emph{Inductive proofs of $q$-log concavity}, Discrete Math.
\textbf{99} (1992), 289--306.

\bibitem{suwangyeh2011}
X.-T.~Su, Y.~Wang, and Y.-N.~Yeh, \emph{Unimodality problems of multinomial
coefficients and symmetric functions}, Electron. J. Combin. \textbf{18(1)}
(2011), \#P73.

\bibitem{burrull-gui-hu}
G.~Burrull, T.~Gui, and H.~Hu, \emph{Asymptotic log-concavity of dominant
lower Bruhat intervals via the Brunn--Minkowski inequality}, arXiv:2311.17980.

\bibitem{carrell1994}
J.~B.~Carrell, \emph{The Bruhat order, algebraic cell decompositions and
rationally smooth Schubert varieties}, Proc. Sympos. Pure Math.
\textbf{56} (1994).

\bibitem{short-intervals}
E.~Akhmedova, \emph{On Bruhat intervals of small lengths for Weyl groups},
arXiv:2110.00862.

\bibitem{petrov1975}
V.~V.~Petrov, \emph{Sums of Independent Random Variables}, Springer, 1975.

\bibitem{equivariant-logconcave}
T.~Gui, \emph{On the equivariant log-concavity for the cohomology of the
flag varieties}, arXiv:2205.05408.

\bibitem{stanley-yan}
S.~H.~Chan and I.~Pak, \emph{Equality cases of the Stanley--Yan log-concave
matroid inequality}, arXiv:2407.19608.

\bibitem{kahn-saks}
R.~van Handel, A.~Yan, and X.~Zeng, \emph{The extremals of the Kahn--Saks
inequality}, arXiv:2309.13434.

\bibitem{gaetzgao2020}
C.~Gaetz and Y.~Gao, \emph{Self-dual intervals in the Bruhat order}, Selecta
Math. (2020). arXiv:2003.06710.

\bibitem{kessouri2024}
K.~Kessouri, M.~Ahmia, A.~Arslan, and A.~Mesbahi, \emph{Combinatorics of
$q$-Mahonian numbers of type B and log-concavity}, arXiv:2408.02424.

\bibitem{chapelier-fromentin}
L.~Chapelier-Laget and G.~Fromentin, \emph{A counterexample to a
Brenti--Carnevale conjecture}, arXiv:2412.19593.

\end{thebibliography}

\end{document}
